\section{Discussion}
\label{sec:discussion}

\subsection{Scope and computational cost}
This paper presents a surrogate-based optimisation framework;
the contribution is the continuous adjoint derivation, SUPG-stabilised
implementation, continuation strategies, and reproducibility infrastructure,
not a directly physically realisable device design.
The Brinkman--SIMP model produces continuously graded permeability
distributions that exploit diffusive mixing through intermediate-density
porous regions; this is a well-established modelling choice in
porous-media topology optimisation~\cite{borrvall2003,lazarov2016}
and is the intended physical model for this class of problems.
The design variable $\rho$ represents local permeability, not a
binary-permeability indicator; the framework optimises within the
Brinkman--SIMP design space.
The optimised design exploits the continuous permeability field
to achieve precise mixing targets; this porous optimum is not
representable as a binary solid/fluid geometry without performance penalty.
\texttt{micrograd} is therefore a tool for surrogate-based inverse
design in porous media, not a direct manufacturable device generator.
Binary realisation, experimental validation, and discrete adjoint
derivation are the principal open directions (Section~\ref{sec:limitations}).

The primary $\nIterPrimary$-iteration optimisation on the $80\times20$ mesh
completes in approximately one hour on a standard laptop (Intel Core i7,
16\,GB RAM); a shorter 600-iteration run completes in under thirty minutes
with RMSE~$=\rmseShortRun$.
Each iteration requires five sparse linear systems (two forward, two
adjoint, one Helmholtz filter), all handled by MUMPS direct factorisation.
For larger meshes the iterative block-preconditioner (Section~\ref{sec:fem})
maintains near-optimal convergence rates (Supplementary~S1).

\subsection{Mesh resolution and optimizer sensitivity}
\label{sec:mesh_limitation}
The $160\times40$ mesh gives RMSE~$=\rmseHiRes$ with a 600-iteration
identical schedule, and RMSE~$=\rmseHiResScaled$ with a resolution-scaled
$\nIterHiRes$-iteration run --- both worse than the $80\times20$ primary
result (RMSE~$=\rmseTopOpt$).
This confirms that the performance difference is not primarily due to
under-iteration: the $160\times40$ optimizer finds a qualitatively
different local minimum with higher gray fraction
($0.557$ vs $0.074$ for the primary mesh).
The topology optimisation problem is multi-modal; the continuation
schedule is calibrated for the $80\times20$ mesh and does not
automatically transfer to finer resolutions.
This study therefore characterises \emph{optimizer sensitivity}
to mesh resolution, not classical Richardson-extrapolation
mesh convergence; the latter would require a mesh-independent
continuation schedule, which is identified as future work.

On coarse meshes ($20\times5$, $N_{\mathrm{out}}=6$) the outlet
discretisation is insufficient to represent the linear gradient target
and the optimizer finds unphysical reversed-flow local minima.
The $40\times10$ mesh ($N_{\mathrm{out}}=11$) converges to a
suboptimal basin; systematic improvement requires the $80\times20$
resolution or finer.
On the primary mesh, the transverse diffusion layer thickness
$\delta \approx 4\,\mu$m (Eq.~\eqref{eq:delta}) is underresolved by
a factor of $\approx 6\times$ ($\Delta y = 25\,\mu$m).
SUPG eliminates spurious oscillations but does \emph{not} recover
sub-grid diffusion structure.
The RMSE metric measures the \emph{global} outlet profile over the full
channel width $L_y = 500\,\mu$m, a scale that is well resolved;
the channel-scale branching topology is therefore correctly captured
even though the wall-normal diffusion sublayer is not.

\subsection{Limitations}
\label{sec:limitations}
\label{sec:known_limitations}
\begin{itemize}
    \item \textbf{Gray-design binary gap.}
          Hard-thresholding the optimised design at $\bar\rho=0.5$
          increases outlet RMSE by $\binaryGapPct\%$
          (Section~\ref{sec:thresholding}), confirming that the surrogate
          exploits gray pathways with no direct binary physical analogue.
          \emph{The Brinkman surrogate RMSE~$=\rmseTopOpt$ is therefore
          not the performance of a physically realisable porous medium};
          it is the optimum within the porous-medium model.
          Systematic binarisation --- through iterative thresholding,
          robust projection~\cite{wang2011}, or post-optimisation
          body-fitted re-solving --- remains an open challenge
          and is deferred to future work.
    \item \textbf{Brinkman approximation.}
          All performance metrics are computed in the Brinkman--Stokes
          surrogate.
          At $\Da \approx 4\times10^{-6} \ll 1$ in solid regions, the
          surrogate accurately represents impermeability~\cite{borrvall2003}.
          Body-fitted Navier--Stokes validation is implemented in
          \texttt{ns\_validation.py} but requires \texttt{pyvista} and
          \texttt{gmsh} dependencies not present in the standard Docker
          image; these results are deferred to a companion study.
    \item \textbf{Two-dimensional steady flow.}
          The framework is 2D; three-dimensional extension is implemented
          in \texttt{validation\_3d.py} but not systematically validated.
    \item \textbf{Fixed objective weights.}
          The weights $w_f$ and $w_c$ are fixed in this study.
          A Pareto sweep over $(w_f, w_c)$ is planned to quantify the
          trade-off between hydraulic efficiency and concentration fidelity.
    \item \textbf{No experimental validation.}
          Results are purely computational.
          Experimental fabrication and testing are necessary next steps.
\end{itemize}



\subsection{Future directions}
Immediate priorities are: (i)~adaptive mesh refinement targeting the
fluid--solid interface; (ii)~integration of body-fitted Navier--Stokes
post-processing into the main pipeline; (iii)~a systematic Pareto front
study over objective weights; (iv)~three-dimensional validation;
and (v)~experimental collaboration for a representative subset of designs.
